\documentclass[11pt]{article}
\usepackage[a4paper,margin=26mm]{geometry}
\usepackage{amsmath,amssymb}
\setlength{\parindent}{0pt}
\setlength{\parskip}{7pt}
\emergencystretch=2em
\title{A nonnormal critical point of a convex spectral function\\
\large Disproof of conjecture 00000007162}
\author{ziangni-sys}
\date{4 October 2026}
\begin{document}
\maketitle
Both source versions assert that critical points of convex spectral
functions are normal matrices. No restriction to symmetric or
Hermitian matrices appears in either version.
On the space of all real $2\times2$ matrices, a nonconstant convex
function of the eigenvalues has a nonnormal critical point.

\textbf{The actual spectral function.}
On $\mathcal M=M_2(\mathbb R)$ define
\[
 f(M)=(\operatorname{tr}M)^2.
\]
Here \emph{spectral} means eigenvalue-dependent, with algebraic
multiplicities. For the two complex eigenvalues $\lambda_1,\lambda_2$
of a real matrix $M$,
\[
 \operatorname{tr}M=\lambda_1+\lambda_2,\qquad
 f(M)=(\lambda_1+\lambda_2)^2.
\]
The sum is real, including for a nonreal conjugate pair.
Thus $f$ is a real-valued spectral function on the entire real
matrix space, not only its symmetric subspace.

Algebraically, trace is the negative coefficient of $t$ in the
actual characteristic polynomial of a $2\times2$ matrix.
Equal characteristic polynomials therefore give equal values of $f$.
Similarity preserves trace and hence $f$.
The function is nonconstant: $f(0)=0$ and $f(I_2)=4$.

\textbf{Convexity on the whole matrix space.}
For $a,b\geq0$ with $a+b=1$ and any $M,N\in\mathcal M$, linearity
of trace gives
\[
 af(M)+bf(N)-f(aM+bN)
 =ab\bigl(\operatorname{tr}M-\operatorname{tr}N\bigr)^2\geq0.
\]
This proves convexity for every pair of real matrices, including
nonnormal matrices.

\newpage
\textbf{A genuine critical point which is not normal.}
For every perturbation $H$, the Fr\'echet derivative is
\[
 Df(M)[H]=2\operatorname{tr}M\,\operatorname{tr}H.
\]
Take the actual Jordan nilpotent
\[
 J=\begin{pmatrix}0&1\\0&0\end{pmatrix}.
\]
Its trace is zero, so $Df(J)$ is the zero functional.
Hence $J$ is a critical point. This is the ordinary Fr\'echet
derivative on a finite-dimensional matrix space; choosing an
equivalent norm does not change the assertion.

A real matrix is normal when $MM^{\mathsf T}=M^{\mathsf T}M$.
But
\[
 JJ^{\mathsf T}=\begin{pmatrix}1&0\\0&0\end{pmatrix},
 \qquad
 J^{\mathsf T}J=\begin{pmatrix}0&0\\0&1\end{pmatrix}.
\]
Their $(0,0)$ entries differ, so $J$ is not normal.

\textbf{Scope.}
This disproves the claimed necessity of normality on unrestricted
real matrices under the eigenvalue-dependent, similarity-invariant
interpretation of spectral function. It does not concern a
singular-value function or a domain pre-restricted to Hermitian
matrices. Within a Hermitian domain normality is automatic.
The source states neither restriction.
The separate minimizer-ordering clause is unnecessary.

\textbf{Formal correspondence.}
Lean constructs the actual real matrices and trace-square function.
\texttt{objective\_spectral} proves dependence on the actual
characteristic polynomial using Mathlib's trace/coefficient theorem.
\texttt{objective\_similarity} proves invariance under every
invertible-matrix conjugation, and \texttt{objective\_convex}
proves \texttt{ConvexOn} on the entire matrix space.

The formal matrix space carries the maximum-entry norm transported
from its finite coordinate space. The actual trace linear map is
converted to a continuous linear map.
\texttt{objective\_derivative} proves \texttt{HasFDerivAt}
with the displayed derivative at every matrix.
\texttt{jordan\_critical} proves zero derivative at $J$, and
\texttt{jordan\_not\_normal} disproves the actual matrix identity.
The final \texttt{counterexample} combines these facts and
nonconstancy.

\textbf{Reproduction.}
Run \texttt{lake build} in \texttt{lean/} with Lean 4.19.0 and
public Mathlib revision
\begin{center}\small\texttt{c44e0c8ee63ca166450922a373c7409c5d26b00b}\end{center}
The project prints the relevant axiom audits.
\end{document}
